
\def\ThesisType{bc}

\def\StudyLanguage{cs}


\def\ThesisTitle{{Optimizing 3D Bin Packing Using Evolutionary Algorithms}}

\def\ThesisAuthor{{Jan Burda}}

\def\YearSubmitted{2026}

\def\Department{{Department of Theoretical Computer Science and Mathematical Logic}}

\def\DeptType{{Department}}

\def\Supervisor{{RNDr. Jiří Fink, Ph.D.}}

\def\SupervisorsDepartment{{Department of Theoretical Computer Science and Mathematical Logic}}

% Study programme (does not apply to rigorosum theses)
\def\StudyProgramme{{Informatika}}

% An optional dedication: you can thank whomever you wish (your supervisor,
% consultant, who provided you with tea and pizza, etc.)
\def\Dedication{%
I would like to thank my supervisor, RNDr. Jiří Fink, Ph.D., 
for his valuable guidance, continuous support, and patience during the course of this work.

\bigskip

\noindent
The ChatGPT tool developed by OpenAI was used to assist in improving the quality 
of the text and code comments. All ideas and conclusions presented in this thesis are original or appropriately referenced.
}

% Abstract (recommended length around 80-200 words; this is not a copy of your thesis assignment!)
\def\Abstract{%
{
The goal of this thesis is to study the three-dimensional bin packing problem with weight 
constraints and to optimize it using nature-inspired algorithms. 
The problem consists of placing a set of three-dimensional boxes into a minimum number of containers with limited capacity.
We design and implement several optimization approaches, including an elitist genetic algorithm and hill-climbing based methods. 
In particular, we investigate different strategies for box ordering and placement heuristic selection, as well as the impact of mutation, 
memetic search, and probabilistic acceptance.
The results show that evolutionary approaches significantly outperform simple local search methods, while certain heuristic 
choices and parameter settings lead to improvements over previously proposed evolutionary approaches.
}
}

% 3 to 5 keywords (recommended) separated by \sep
% Keywords are useful for indexing and searching for the theses by topic.
\def\ThesisKeywords{%
3D Bin Packing \sep Evolutionary Algorithms \sep Genetic Algorithms \sep Local Search
}

% If any of your metadata strings contains TeX macros, you need to provide
% a plain-text version for use in XMP metadata embedded in the output PDF file.
% If you are not sure, check the generated thesis.xmpdata file.
\def\ThesisAuthorXMP{\ThesisAuthor}
\def\ThesisTitleXMP{\ThesisTitle}
\def\ThesisKeywordsXMP{\ThesisKeywords}
\def\AbstractXMP{\Abstract}

% If your abstracts are long and do not fit in the infopage, you can make the
% fonts a bit smaller by this setting. (Also, you should try to compress your abstract more.)
\def\InfoPageFont{}
%\def\InfoPageFont{\small}  % uncomment to decrease font size



\def\ThesisTitleCS{{Optimalizace problému 3D pakování pomocí evolučních algoritmů}}
\def\DepartmentCS{{Katedra teoretické informatiky a matematické logiky}}
\def\DeptTypeCS{{Katedra}}
\def\SupervisorsDepartmentCS{{Katedra teoretické informatiky a matematické logiky}}
\def\StudyProgrammeCS{{Informatika}}

\def\ThesisKeywordsCS{%
3D pakování \sep Evoluční agloritmy \sep Genetické algoritmy \sep Lokální prohledávání
}

\def\AbstractCS{%
{
Cílem této práce je studium problému trojrozměrného bin packingu s omezením na hmotnost 
a jeho optimalizace pomocí přírodou inspirovaných algoritmů. 
Úloha spočívá v umístění množiny trojrozměrných kvádrů do minimálního počtu kontejnerů s omezenou kapacitou.
Navrhujeme a implementujeme několik optimalizačních přístupů, včetně elitistického genetického algoritmu 
a metod založených na hill climbingu. 
Zaměřujeme se zejména na různé strategie uspořádání kvádrů a volby heuristik pro jejich umísťování, 
stejně jako na vliv mutace, memetického prohledávání a pravděpodobnostního přijímání horších řešení.
Výsledky ukazují, že evoluční přístupy výrazně překonávají jednoduché metody lokálního prohledávání, 
zatímco vhodná volba heuristik a nastavení parametrů vede ke zlepšení oproti dříve navrženým evolučním přístupům.
}
}
